\makeatletter
\@ifundefined{keepfig}{\newenvironment{keepfig}{\par\edef\keepfig@pi{\the\parindent}\noindent
  \begin{minipage}{\linewidth}\parindent=\keepfig@pi\relax}{\end{minipage}\par}}{}
\makeatother
\section{Круговые расширения}

Элемент $\zeta$ поля $F$ называется \emph{первообразным корнем $n$-й степени из $1$},
если его порядок в группе $F^\times$ равен $n$. Корни $n$-й степени из $1$ в $F$ образуют подгруппу
$\mu_n(F)=\{x\in F\mid x^n=1\}$ группы $F^\times$; в ней не больше $n$ элементов, так как у $X^n-1$ не больше $n$ корней.

Первообразный корень $n$-й степени из $1$ может существовать, лишь если характеристика $F$ не делит $n$ (т.\,е.\ равна
$0$ или простому $p\nmid n$). Действительно, если $\operatorname{char}F=p$ и $n=pm$, то $X^n-1=(X^m-1)^p$ (в
характеристике $p$ возведение в степень $p$~--- гомоморфизм), поэтому из $x^n=1$ следует $(x^m-1)^p=0$, т.\,е.\
$x^m=1$, и элементов порядка $n$ нет.

\subsection*{Корни из единицы и циклические группы}

Нам понадобятся два простых факта о циклических группах и о корнях из единицы. Через $\varphi(n)$ обозначается
функция Эйлера~--- число целых $1\le i\le n$, взаимно простых с $n$; это порядок группы $(\Z/n\Z)^\times$.

\begin{lemma}[5.7b]
Пусть $C=\langle\zeta\rangle$~--- циклическая группа порядка $n$.
\begin{enumerate}[label=\textup{(\alph*)}]
\item Порядок элемента $\zeta^i$ равен $n/\gcd(i,n)$. В частности, $\zeta^i$ порождает $C$ тогда и только тогда, когда
$\gcd(i,n)=1$, и образующих ровно $\varphi(n)$.
\item Для каждого $d\mid n$ в $C$ есть ровно одна подгруппа порядка $d$, а именно $\{x\in C\mid x^d=1\}=\langle\zeta^{n/d}\rangle$,
и ровно $\varphi(d)$ элементов порядка $d$. Поэтому $\sum_{d\mid n}\varphi(d)=n$.
\end{enumerate}
\end{lemma}

\begin{proof}
(a) Пусть $e=\gcd(i,n)$. Тогда $(\zeta^i)^k=1\iff n\mid ik\iff\frac ne\mid\frac ie\,k\iff\frac ne\mid k$, так как
$\gcd(\frac ne,\frac ie)=1$. (b) Элемент $\zeta^j$ удовлетворяет $(\zeta^j)^d=1$ тогда и только тогда, когда $n\mid jd$,
т.\,е.\ $\frac nd\mid j$; таких классов $j\bmod n$ ровно $d$, и они образуют подгруппу $\langle\zeta^{n/d}\rangle$. Всякая
подгруппа порядка $d$ состоит из элементов с $x^d=1$ (теорема Лагранжа) и потому совпадает с ней. Элементы порядка $d$
лежат в этой циклической подгруппе порядка $d$ и являются её образующими; по (a) их $\varphi(d)$. Подсчитав элементы $C$
по их порядкам, получаем $\sum_{d\mid n}\varphi(d)=n$.
\end{proof}

В частности, \emph{если $\zeta\in F$~--- первообразный корень $n$-й степени из $1$, то $\mu_n(F)=\langle\zeta\rangle$
(в $\mu_n(F)$ не больше $n$ элементов, а $n$ различных степеней $\zeta$ в нём уже есть), и первообразные корни $n$-й степени
из $1$ в $F$~--- это в точности элементы $\zeta^i$ с $\gcd(i,n)=1$; их $\varphi(n)$}: это элементы порядка $n$ циклической
группы $\langle\zeta\rangle$.

\begin{lemma}[5.7c]
Каждая конечная подгруппа $G$ мультипликативной группы $F^\times$ поля $F$ циклическая.
\end{lemma}

\begin{proof}
Пусть $|G|=n$ и $\psi(d)$~--- число элементов порядка $d$ в $G$ ($d\mid n$). Если $\psi(d)\ne0$ и $x\in G$ имеет порядок
$d$, то все $d$ элементов $\langle x\rangle$~--- корни многочлена $X^d-1$, у которого в поле не больше $d$ корней.
Поэтому все $y\in G$ с $y^d=1$ лежат в $\langle x\rangle$, в частности элементы порядка $d$~--- это образующие
$\langle x\rangle$, и $\psi(d)=\varphi(d)$ по \rf{5.7b}. Итак, $\psi(d)\le\varphi(d)$ для всех $d\mid n$, и
\[
n=\sum_{d\mid n}\psi(d)\le\sum_{d\mid n}\varphi(d)=n .
\]
Значит, всюду равенство, $\psi(n)=\varphi(n)\ge1$, и в $G$ есть элемент порядка $n$.
\end{proof}

Существенно, что $F$~--- поле: в кольце $\Z/8\Z$ у уравнения $x^2=1$ четыре решения $1,3,5,7$, и группа
$(\Z/8\Z)^\times=\{1,3,5,7\}$ не циклическая.

\subsection*{Поле разложения многочлена $X^n-1$}

\begin{proposition}[5.8]
Пусть характеристика поля $F$ не делит $n$, и $E$~--- поле разложения многочлена $X^n-1$ над $F$.
\begin{enumerate}[label=\textup{(\alph*)}]
\item В $E$ есть первообразный корень $n$-й степени из $1$.
\item Если $\zeta\in E$~--- первообразный корень $n$-й степени из $1$, то $E=F[\zeta]$.
\item Расширение $E/F$~--- расширение Галуа. Для каждого $\sigma\in\Gal(E/F)$ существует $i\in(\Z/n\Z)^\times$ такое, что
$\sigma\zeta=\zeta^i$ для \emph{всех} $\zeta$ с $\zeta^n=1$; отображение $\sigma\mapsto[i]$~--- инъективный гомоморфизм
\[
\Gal(E/F)\to(\Z/n\Z)^\times .
\]
\end{enumerate}
\end{proposition}

\begin{proof}
(a) Производная многочлена $X^n-1$ равна $nX^{n-1}$, причём $n\ne0$ в $F$ (здесь используется условие на
характеристику). Её корни (если они есть) равны $0$, а $0$ не является корнем $X^n-1$, поэтому по \rf{2.19} у $X^n-1$ нет
кратных корней, и
в $E$ лежат $n$ различных корней $n$-й степени из $1$. Они образуют подгруппу $\mu_n(E)$ порядка $n$ в $E^\times$,
циклическую по \rf{5.7c}. Её образующая имеет порядок $n$, т.\,е.\ является первообразным корнем.

(b) Корни $X^n-1$~--- это $n$ различных степеней $1,\zeta,\dots,\zeta^{n-1}$, т.\,е.\
$X^n-1=\prod_{k=0}^{n-1}(X-\zeta^k)$, и все они лежат в $F[\zeta]$. Значит, $E=F[\zeta]$.

(c) Поле $E$~--- поле разложения сепарабельного многочлена, поэтому $E/F$~--- расширение Галуа (\rf{3.10}). Пусть
$\sigma\in\Gal(E/F)$ и $\zeta$~--- первообразный корень. Ограничение $\sigma$ на $E^\times$~--- инъективный гомоморфизм
групп, поэтому $\sigma\zeta$ тоже имеет порядок $n$, и по \rf{5.7b} $\sigma\zeta=\zeta^i$ с $\gcd(i,n)=1$; класс $[i]$ определён
однозначно. Каждый корень $n$-й степени из $1$ имеет вид $\xi=\zeta^m$, и
\[
\sigma\xi=(\sigma\zeta)^m=\zeta^{im}=\xi^i,
\]
так что одно и то же $i$ годится для всех корней $n$-й степени из $1$ (в частности, $[i]$ не зависит от выбора
первообразного корня $\zeta$). Если $\sigma\zeta=\zeta^i$ и $\tau\zeta=\zeta^j$, то
$\sigma\tau(\zeta)=\sigma(\zeta^j)=\zeta^{ij}$, т.\,е.\ отображение~--- гомоморфизм. Если $[i]=1$, то $\sigma\zeta=\zeta$, и
$\sigma=\id$, так как $\zeta$ порождает $E$ над $F$; значит, гомоморфизм инъективен.
\end{proof}

\begin{remark}[5.8a: сложение и умножение в $\Z/n\Z$]
Зафиксируем первообразный корень $\zeta$. Отображение $a\mapsto\zeta^a$~--- изоморфизм аддитивной группы $(\Z/n\Z,+)$ на
$(\mu_n,\cdot)$: \emph{умножению корней отвечает сложение показателей}, и порядок $\zeta^a$ равен аддитивному порядку
$n/\gcd(a,n)$ класса $a$. Возведение в степень $i$ умножает показатели: $(\zeta^a)^i=\zeta^{ai}$. Каждый автоморфизм
циклической группы $\mu_n$ имеет такой вид, так что $\operatorname{Aut}(\mu_n)\cong(\Z/n\Z)^\times$, и гомоморфизм \rf{5.8}(c)~---
это просто действие группы Галуа на $\mu_n$:
\[
\Gal(E/F)\to\operatorname{Aut}(\mu_n)=(\Z/n\Z)^\times .
\]
Так в теории круговых полей работают обе операции кольца $\Z/n\Z$: сложение описывает саму группу $\mu_n$, а умножение~---
действие Галуа на ней. Первообразные корни отвечают подмножеству $(\Z/n\Z)^\times\subset\Z/n\Z$ обратимых показателей.
\end{remark}

\begin{example}[5.8b]
Гомоморфизм $\Gal(F[\zeta]/F)\to(\Z/n\Z)^\times$ не обязан быть сюръективным.
\begin{itemize}[leftmargin=*]
\item $F=\CC$: здесь $E=\CC$, и образ равен $\{1\}$.
\item $F=\R$: при $n\le2$ все корни вещественны, $E=\R$ и образ $\{[1]\}$; при $n\ge3$ имеем $\zeta\notin\R$, $E=\CC$, и
образ $\{[1],[-1]\}$: нетривиальный автоморфизм~--- комплексное сопряжение, а на корнях из единицы (они лежат на
единичной окружности) $\bar\zeta=\zeta^{-1}$.
\item $F=\Q$, $n=p$ простое: $[\Q[\zeta]:\Q]=p-1$ по \rf{1.42}, и гомоморфизм биективен. Ниже (\rf{5.10}) мы докажем, что над $\Q$ он
биективен при всех $n$.
\item $F=\F_p$, $p\nmid n$: образ~--- циклическая подгруппа, порождённая классом $[p]$ (предложение \rf{5.8c}).
\end{itemize}
\begin{keepfig}
\begin{center}
\begin{tikzpicture}[scale=1.15,>=stealth]
\foreach \n/\sh in {5/0,6/4.6}{
\begin{scope}[xshift=\sh cm]
\draw[->,gray] (-1.45,0) -- (1.5,0);
\draw[->,gray] (0,-1.35) -- (0,1.4);
\draw[cold] (0,0) circle (1);
\foreach \k in {1,...,\n}{\fill[cold] ({cos(360*\k/\n)},{sin(360*\k/\n)}) circle (1.6pt);}
\draw[<->,cnew,dashed,shorten <=3pt,shorten >=3pt] ({cos(360/\n)},{sin(360/\n)}) -- ({cos(360/\n)},{-sin(360/\n)});
\node[cold,above right,font=\footnotesize] at ({cos(360/\n)},{sin(360/\n)}) {$\zeta$};
\node[cold,below right,font=\footnotesize] at ({cos(360/\n)},{-sin(360/\n)}) {$\bar\zeta=\zeta^{-1}$};
\node[font=\footnotesize] at (0,-1.7) {$n=\n$};
\end{scope}}
\node[above right,font=\footnotesize] at (1,0) {$1$};
\node[above right,font=\footnotesize] at (5.6,0) {$1$};
\node[above left,font=\footnotesize] at (3.6,0) {$-1$};
\end{tikzpicture}
\end{center}
\figcap{Корни степени $5$ и $6$ из $1$ на единичной окружности. Комплексное сопряжение~--- отражение относительно
вещественной оси~--- переводит $\zeta$ в $\zeta^{-1}$, т.\,е.\ действует как $[-1]\in(\Z/n\Z)^\times$; на месте остаются
лишь вещественные корни $1$ (и $-1$ при чётном $n$).}
\end{keepfig}
\end{example}

\begin{proposition}[5.8c: корни из единицы над $\F_p$]
Пусть $p$~--- простое число, $m=p^kn$, где $p\nmid n$, и $f$~--- порядок класса $[p]$ в $(\Z/n\Z)^\times$ (наименьшее $f\ge1$
с $p^f\equiv1\pmod n$). Тогда поле разложения многочлена $X^m-1$ над $\F_p$ есть $\F_{p^f}$, а образ его группы Галуа в
$(\Z/n\Z)^\times$ при гомоморфизме \rf{5.8}(c)~--- подгруппа $\langle[p]\rangle$.
\end{proposition}

\begin{proof}
В характеристике $p$ имеем $X^m-1=(X^n)^{p^k}-1=(X^n-1)^{p^k}$, так что корни у $X^m-1$ те же, что у $X^n-1$ (только
с кратностью $p^k$), и поля разложения совпадают. Пусть $E=\F_p[\zeta]$~--- поле разложения $X^n-1$ (\rf{5.8}). Группа
$\Gal(E/\F_p)$ порождается автоморфизмом Фробениуса $x\mapsto x^p$ (\rf{4.20}), а он переводит $\zeta$ в $\zeta^p$, т.\,е.\
ему отвечает класс $[p]$. Значит, образ группы Галуа~--- $\langle[p]\rangle$, подгруппа порядка $f$, и по инъективности
$[E:\F_p]=f$, т.\,е.\ $E\cong\F_{p^f}$.
\end{proof}

То же можно увидеть и иначе: $\F_{p^e}^\times$~--- циклическая группа порядка $p^e-1$, и по \rf{5.7b}(b) в ней есть элемент
порядка $n$ (и тогда подгруппа порядка $n$ в ней~--- в точности $\{x\mid x^n=1\}=\mu_n$) тогда и только тогда, когда
$n\mid p^e-1$; наименьшее такое $e$ равно $f$. Например:
\begin{itemize}[leftmargin=*]
\item $X^7-1$ над $\F_2$: $2^3=8\equiv1\pmod7$, $f=3$, поле разложения $\F_8$, и $X^7-1=(X+1)(X^3+X+1)(X^3+X^2+1)$;
\item $X^{12}-1$ над $\F_3$: $12=3\cdot4$, $X^{12}-1=(X^4-1)^3=(X-1)^3(X+1)^3(X^2+1)^3$, $f=\operatorname{ord}_4(3)=2$, поле
разложения $\F_9$;
\item $X^{p-1}-1$ раскладывается уже над $\F_p$ ($f=1$)~--- это малая теорема Ферма; $X^{p^k}-1=(X-1)^{p^k}$.
\end{itemize}

\subsection*{Круговые многочлены}

Отображение $\sigma\mapsto[i]$ из \rf{5.8} не всегда сюръективно, но над $\Q$ оно сюръективно при всех $n$. Чтобы это
доказать, введём многочлены, корни которых~--- в точности первообразные корни из $1$.

\begin{definition}
\emph{$n$-й круговой многочлен}~--- это
\[
\Phi_n(X)=\prod(X-\zeta)\qquad\text{(произведение по всем первообразным корням $n$-й степени из $1$ в $\CC$)}.
\]
\end{definition}

Его степень равна $\varphi(n)$ (\rf{5.7b}). Каждый корень $n$-й степени из $1$ имеет некоторый порядок $d\mid n$ и является
первообразным корнем степени $d$ ровно для одного $d$; поэтому
\begin{equation}\label{eq:W1-Xn-1}
X^n-1=\prod_{d\mid n}\Phi_d(X).
\end{equation}
Отсюда $\Phi_n=(X^n-1)\big/\prod_{d\mid n,\,d<n}\Phi_d$~--- удобный индуктивный способ вычисления. Иначе говоря, $\Phi_n$
получается из $X^n-1$ удалением корней всех многочленов $X^d-1$ с $d\mid n$, $d<n$, \emph{каждого по одному разу}
(делить на произведение всех $X^d-1$ нельзя: при $n=6$ произведение $(X-1)(X^2-1)(X^3-1)$ имеет ту же степень $6$).
Например, $\Phi_1=X-1$, $\Phi_2=X+1$, $\Phi_3=X^2+X+1$ и
\[
\Phi_6=\frac{X^6-1}{(X-1)(X+1)(X^2+X+1)}=\frac{X^6-1}{(X+1)(X^3-1)}=\frac{X^3+1}{X+1}=X^2-X+1 .
\]
Ещё несколько:
\[
\Phi_4=X^2+1,\quad \Phi_8=X^4+1,\quad \Phi_9=X^6+X^3+1,\quad \Phi_{10}=X^4-X^3+X^2-X+1,\quad \Phi_{12}=X^4-X^2+1 .
\]
Общие правила: $\Phi_p=X^{p-1}+\dots+X+1$ для простого $p$ (так как $X^p-1=(X-1)\Phi_p$); $\Phi_{p^k}(X)=\Phi_p(X^{p^{k-1}})$
(первообразные корни степени $p^k$~--- это корни $X^{p^k}-1$, не являющиеся корнями $X^{p^{k-1}}-1$, а
$(X^{p^k}-1)/(X^{p^{k-1}}-1)=\Phi_p(X^{p^{k-1}})$); $\Phi_{2m}(X)=\Phi_m(-X)$ для нечётного $m>1$ (при нечётном $m$
отображение $\zeta\mapsto-\zeta$ инъективно переводит первообразные корни степени $m$ в первообразные корни степени $2m$,
а их поровну: $\varphi(2m)=\varphi(m)$; кроме того, степень $\varphi(m)$ чётна, так что
$\Phi_{2m}(X)=\prod(X+\zeta)=(-1)^{\varphi(m)}\prod(-X-\zeta)=\Phi_m(-X)$). У $\Phi_n$ с $n<105$ все коэффициенты
равны $0$ или $\pm1$, но у $\Phi_{105}$ (степени $48$) коэффициенты при $X^7$ и $X^{41}$ равны $-2$ (это проверяется
прямым вычислением).

\begin{proposition}[5.8d]
\begin{enumerate}[label=\textup{(\alph*)}]
\item $\Phi_n$~--- унитарный многочлен с целыми коэффициентами.
\item Пусть $F$~--- поле, характеристика которого не делит $n$, и $E$~--- поле разложения $X^n-1$ над $F$. Тогда образ
$\Phi_n$ в $F[X]$ равен $\prod(X-\xi)$, где $\xi$ пробегает первообразные корни $n$-й степени из $1$ в $E$; в частности, его
корни~--- это в точности первообразные корни $n$-й степени из $1$, и все они простые.
\end{enumerate}
\end{proposition}

\begin{proof}
(a) Индукция по $n$. Произведение $\prod_{d\mid n,\,d<n}\Phi_d$ унитарно и по предположению индукции лежит в $\Z[X]$; при
делении с остатком на унитарный многочлен с целыми коэффициентами частное и остаток тоже целые, а по \eqref{eq:W1-Xn-1}
частное от деления $X^n-1$ равно $\Phi_n$. (Можно сослаться и на \rf{1.14}: унитарный делитель в $\Q[X]$ унитарного многочлена из
$\Z[X]$ лежит в $\Z[X]$.)

(b) Равенство \eqref{eq:W1-Xn-1} верно в $\Z[X]$ и потому в $F[X]$ (для образов). Для $d\mid n$ положим
$\Psi_d=\prod(X-\xi)\in E[X]$, где $\xi$ пробегает элементы порядка $d$ группы $\mu_n(E)$. Группа $\mu_n(E)$ циклическая
порядка $n$ (\rf{5.8}), её подгруппа $\mu_d(E)$ имеет порядок $d$ (\rf{5.7b}), и, разбивая её элементы по порядкам, получаем
$X^d-1=\prod_{e\mid d}\Psi_e$ в $E[X]$. Индукцией по $d\mid n$ получаем $\Phi_d=\Psi_d$ в $E[X]$: при $d=1$ оба равны $X-1$,
а если это верно для всех собственных делителей $d$, то сокращаем равные множители в
$\prod_{e\mid d}\Phi_e=X^d-1=\prod_{e\mid d}\Psi_e$ (в $E[X]$ нет делителей нуля). Простота корней следует из
сепарабельности $X^n-1$ (доказательство \rf{5.8}(a)).
\end{proof}

\begin{lemma}[5.9]
Пусть характеристика $F$ не делит $n$, и $\zeta$~--- первообразный корень $n$-й степени из $1$ в некотором расширении $F$.
Следующие условия равносильны:
\begin{enumerate}[label=\textup{(\alph*)}]
\item круговой многочлен $\Phi_n$ неприводим в $F[X]$;
\item $[F[\zeta]:F]=\varphi(n)$;
\item гомоморфизм $\Gal(F[\zeta]/F)\to(\Z/n\Z)^\times$ из \rf{5.8} является изоморфизмом.
\end{enumerate}
\end{lemma}

\begin{proof}
По \rf{5.8d}(b) $\zeta$~--- корень $\Phi_n$, поэтому минимальный многочлен $g$ элемента $\zeta$ над $F$ делит $\Phi_n$, и
$[F[\zeta]:F]=\deg g$. (a)$\Rightarrow$(b): если $\Phi_n$ неприводим, то $g=\Phi_n$ (оба унитарны), и $\deg g=\varphi(n)$.
(b)$\Rightarrow$(a): если $\deg g=\varphi(n)=\deg\Phi_n$, то $g=\Phi_n$, и $\Phi_n$ неприводим. (b)$\iff$(c): $F[\zeta]$~---
поле разложения $X^n-1$ (\rf{5.8}(b)), поэтому $|\Gal(F[\zeta]/F)|=[F[\zeta]:F]$, а инъективный гомоморфизм в группу порядка
$\varphi(n)$ биективен ровно тогда, когда $|\Gal(F[\zeta]/F)|=\varphi(n)$.
\end{proof}

\begin{theorem}[5.10]
Круговой многочлен $\Phi_n$ неприводим в $\Q[X]$.
\end{theorem}

\begin{proof}[Первое доказательство: редукция по модулю $p$]
Пусть $f$~--- унитарный неприводимый делитель $\Phi_n$ в $\Q[X]$; по \rf{1.14} $f\in\Z[X]$. Его корни~--- некоторые
первообразные корни $n$-й степени из $1$, и нужно показать, что это все первообразные корни. Так как все они имеют вид
$\zeta^i$ с $\gcd(i,n)=1$ (\rf{5.7b}), достаточно доказать:
\[
\zeta\ \text{корень}\ f\ \Longrightarrow\ \zeta^i\ \text{корень}\ f\quad\text{для всех $i\ge1$ с}\ \gcd(i,n)=1 .
\]
Такое $i$~--- произведение простых $p_1p_2\cdots p_r$ (с повторениями), ни одно из которых не делит $n$. Поэтому
достаточно доказать
\[
(\ast)\qquad \zeta\ \text{корень}\ f,\ \ p\ \text{простое},\ p\nmid n\ \Longrightarrow\ \zeta^p\ \text{корень}\ f:
\]
применяя $(\ast)$ последовательно к корням $\zeta$, $\zeta^{p_1}$, $\zeta^{p_1p_2}$, \dots, получим, что $\zeta^i$~--- корень $f$.

Докажем $(\ast)$. Запишем $\Phi_n=fg$; по \rf{1.14} $g\in\Z[X]$, $g$ унитарен. Предположим, что $\zeta^p$ не корень $f$. Так
как $\gcd(p,n)=1$, $\zeta^p$~--- первообразный корень и, значит, корень $\Phi_n$; поэтому $g(\zeta^p)=0$, т.\,е.\ $\zeta$~---
корень многочлена $g(X^p)$. Итак, неприводимый $f$ и $g(X^p)$ имеют общий корень, и по \rf{2.17} $f\mid g(X^p)$ в $\Q[X]$:
$g(X^p)=f\,h$, где $h\in\Z[X]$ (частное унитарных многочленов из $\Z[X]$).

Обозначим через $u\mapsto\bar u$ редукцию $\Z[X]\to\F_p[X]$. Так как $a^p=a$ для $a\in\F_p$, а возведение в степень $p$
в характеристике $p$ аддитивно,
\[
\bar g(X)^p=\Bigl(\sum\bar a_iX^i\Bigr)^p=\sum\bar a_i^{\,p}X^{ip}=\sum\bar a_iX^{ip}=\bar g(X^p).
\]
Поэтому $\bar f\mid\bar g^{\,p}$ в $\F_p[X]$. Многочлен $\bar f$ унитарен степени $\deg f\ge1$; пусть $q$~--- его неприводимый
множитель. Тогда $q\mid\bar g^{\,p}$, а значит, $q\mid\bar g$, и $q^2\mid\bar f\bar g=\bar\Phi_n$, а $\bar\Phi_n$ делит
$X^n-1$ в $\F_p[X]$ по \eqref{eq:W1-Xn-1}. Но при $p\nmid n$ многочлен $X^n-1$ сепарабелен над $\F_p$: корни его
производной $nX^{n-1}$ (где $n\ne0$ в $\F_p$), если они есть, равны $0$, а $0$ не корень $X^n-1$. А многочлен, делящийся на $q^2$, имеет
кратный корень (любой корень $q$)~--- противоречие.
\end{proof}

\begin{proof}[Второе доказательство: элементы Фробениуса]
По \rf{5.9} нужно доказать, что гомоморфизм $\rho\colon\Gal(\Q[\zeta]/\Q)\to(\Z/n\Z)^\times$ сюръективен. Пусть $p\nmid n$~---
простое число. Применим конструкцию из доказательства теоремы Дедекинда \rf{4.28} к многочлену $X^n-1\in\Z[X]$: он унитарен, и
его редукция по модулю $p$ сепарабельна. Его корни $1,\zeta,\dots,\zeta^{n-1}$ порождают кольцо $A=\Z[\zeta]$; пусть
$P$~--- максимальный идеал $A$, содержащий $p$, и $P\cap\Z=p\Z$. Такой $P$ существует: $A$~--- конечно порождённая
абелева группа (её порождают $1,\zeta,\dots,\zeta^{n-1}$) без кручения, т.\,е.\ $A\cong\Z^r$, $r\ge1$, поэтому
$A/pA\cong\F_p^r\ne0$; в конечном ненулевом кольце $A/pA$ есть максимальный идеал, и его прообраз в $A$~--- искомый $P$
(а $P\cap\Z$~--- собственный идеал $\Z$, содержащий $p$, т.\,е.\ $p\Z$). Доказательство \rf{4.28} даёт элемент $\sigma_P\in G=\Gal(\Q[\zeta]/\Q)$ с
\[
\sigma_P(a)\equiv a^p\pmod P\qquad\text{для всех } a\in A
\]
(\emph{элемент Фробениуса} в $P$, \rf{4.29}). В частности, $\sigma_P(\zeta)\equiv\zeta^p\pmod P$. С другой стороны,
$\sigma_P(\zeta)=\zeta^j$ для некоторого $j$. Корни $n$-й степени из $1$ попарно несравнимы по модулю $P$: подставив $X=1$
в равенство $X^{n-1}+\dots+X+1=\prod_{k=1}^{n-1}(X-\zeta^k)$, получаем
\[
\prod_{k=1}^{n-1}(1-\zeta^k)=n,
\]
а $n\notin P$ (иначе $n\in P\cap\Z=p\Z$); так как идеал $P$ простой, $1-\zeta^k\notin P$ при $1\le k\le n-1$. Если бы
$\zeta^a\equiv\zeta^b\pmod P$ при $a\not\equiv b\pmod n$, то $\zeta^a(1-\zeta^{b-a})\in P$, и, поскольку $\zeta^a$
обратим в $A$, $1-\zeta^{b-a}\in P$~--- противоречие. Значит, $\zeta^j=\zeta^p$, т.\,е.\ $\sigma_P(\zeta)=\zeta^p$ и
$\rho(\sigma_P)=[p]$.

Осталось заметить, что классы простых $p\nmid n$ порождают $(\Z/n\Z)^\times$. Каждый элемент $(\Z/n\Z)^\times$ имеет вид
$[a]$ с $1\le a\le n$, $\gcd(a,n)=1$; разложим $a=p_1\cdots p_r$ на простые множители. Ни одно $p_i$ не делит $n$ (иначе
$p_i\mid\gcd(a,n)$), поэтому $[a]=[p_1]\cdots[p_r]=\rho(\sigma_{P_1}\cdots\sigma_{P_r})$, где $P_i$~--- максимальный идеал над
$p_i$. Итак, $\rho$ сюръективен, и $\Phi_n$ неприводим по \rf{5.9}.
\end{proof}

Во втором доказательстве используются лишь простые $p\nmid n$ и лишь то, что каждый класс из $(\Z/n\Z)^\times$~---
\emph{произведение} классов таких простых; теорема Дирихле о том, что в каждом классе есть простое число, не нужна. Роль
элемента Фробениуса здесь такова: он даёт для каждого $p\nmid n$ автоморфизм $\sigma_P$, действие которого на $\zeta$
известно заранее ($\zeta\mapsto\zeta^p$), хотя мы ещё не знаем, что $\zeta^p$ сопряжён с $\zeta$; этот автоморфизм получается
подъёмом в $G$ автоморфизма Фробениуса $x\mapsto x^p$ поля $A/P$ (\rf{4.28}). В обоих доказательствах решающую роль играет
сепарабельность $X^n-1$ по модулю $p$: она позволяет различать корни из единицы после редукции.

\begin{aside}[5.11]
Доказательство \rf{5.10} очень старое~--- по существу оно восходит к Дедекинду (1857),~--- но его схема стала популярной
сравнительно недавно: утверждение в характеристике $0$ редуцируют по модулю $p$ (где оно может перестать быть верным) и
используют существование автоморфизма Фробениуса $a\mapsto a^p$, чтобы доказать исходное утверждение. Так доказывают,
например, некоторые теоремы о коммутативных кольцах и о комплексных многообразиях. Есть и глубокие связи между
характеристиками $0$ и $p$: для $f\in\Z[X_1,\dots,X_n]$ можно рассмотреть (a) топологическое пространство решений $f=0$ в
$\CC^n$ и (b) множества решений $\bar f=0$ в конечных полях $\F_{p^m}$; гипотезы Вейля (доказанные Дворком, Гротендиком
и Делинем) утверждают, что
числа Бетти пространства (a) управляют количествами точек в (b).
\end{aside}

\subsection*{Общая теория круговых полей $\Q(\zeta_n)$}

Далее $\zeta_n=e^{2\pi i/n}\in\CC$. Поле $\Q(\zeta_n)=\Q[\zeta_n]$ не зависит от выбора первообразного корня (все они~---
степени друг друга) и называется \emph{$n$-м круговым полем}; $\Q(\zeta_1)=\Q(\zeta_2)=\Q$. Для $a\in(\Z/n\Z)^\times$ через
$\sigma_a$ обозначим автоморфизм $\zeta\mapsto\zeta^a$ ($\zeta\in\mu_n$).

\begin{theorem}[5.11a]
\begin{enumerate}[label=\textup{(\alph*)}]
\item $\Q(\zeta_n)/\Q$~--- расширение Галуа степени $\varphi(n)$, а $a\mapsto\sigma_a$~--- изоморфизм
$(\Z/n\Z)^\times\xrightarrow{\ \sim\ }\Gal(\Q(\zeta_n)/\Q)$; в частности, группа Галуа абелева.
\item Подполя $K\subset\Q(\zeta_n)$ взаимно однозначно соответствуют подгруппам $H\subset(\Z/n\Z)^\times$:
$K=\Q(\zeta_n)^H$, $H=\Gal(\Q(\zeta_n)/K)$, $[K:\Q]=\varphi(n)/|H|$. Каждое подполе~--- расширение Галуа поля $\Q$ с группой
$(\Z/n\Z)^\times/H$.
\end{enumerate}
\end{theorem}

\begin{proof}
(a) Это \rf{5.8}, \rf{5.9} и \rf{5.10}. (b) Основная теорема теории Галуа (\rf{3.17}); подгруппы абелевой группы нормальны, поэтому все
подполя~--- расширения Галуа, и $\Gal(K/\Q)\cong(\Z/n\Z)^\times/H$.
\end{proof}

Так описание подполей кругового поля сводится к задаче о подгруппах конечной абелевой группы $(\Z/n\Z)^\times$. Её
строение известно из элементарной теории чисел (приводим его без доказательства; в этой главе оно не используется): по
китайской теореме об остатках $(\Z/n\Z)^\times\cong\prod_{p\mid n}(\Z/p^{k_p}\Z)^\times$
($n=\prod p^{k_p}$); группа $(\Z/p^k\Z)^\times$ циклическая порядка $p^{k-1}(p-1)$ при нечётном $p$, а
$(\Z/2^k\Z)^\times\cong\Z/2\times\Z/2^{k-2}$ при $k\ge2$.

\begin{proposition}[5.11b: корни из единицы в $\Q(\zeta_n)$]
\begin{enumerate}[label=\textup{(\alph*)}]
\item Если $n$ нечётно, то $-\zeta_n$~--- первообразный корень степени $2n$ из $1$, и $\Q(\zeta_{2n})=\Q(\zeta_n)$.
\item Корни из единицы в $\Q(\zeta_n)$~--- это $\mu_n$ при чётном $n$ и $\mu_{2n}$ при нечётном $n$.
\item $\Q(\zeta_m)\subset\Q(\zeta_n)$ тогда и только тогда, когда $m\mid n$ либо $n$ нечётно и $m\mid2n$. В частности,
$\Q(\zeta_m)=\Q(\zeta_n)$ лишь при $m=n$ или $\{m,n\}=\{k,2k\}$ с нечётным $k$; поэтому обычно считают $n\not\equiv2\pmod4$.
\end{enumerate}
\end{proposition}

\begin{proof}
(a) Элементы $-1$ и $\zeta_n$ перестановочны и имеют взаимно простые порядки $2$ и $n$, поэтому порядок $-\zeta_n$ равен
$2n$. Так как $\Q(-\zeta_n)=\Q(\zeta_n)$, получаем $\Q(\zeta_{2n})=\Q(\zeta_n)$.

(b) Пусть $\xi\in\Q(\zeta_n)$~--- корень из $1$ порядка $m$. Подгруппа $\langle\xi,\zeta_n\rangle\subset\Q(\zeta_n)^\times$
конечна (порождена двумя перестановочными элементами конечного порядка), поэтому циклическая (\rf{5.7c}); её порядок $L$
делится на $m$ и на $n$, а её образующая $\eta$~--- первообразный корень степени $L$ из $1$, лежащий в $\Q(\zeta_n)$. Значит,
$\varphi(L)=[\Q(\eta):\Q]\le\varphi(n)$ и $n\mid L$. Пусть $L=nt$. Так как $\varphi(x)=x\prod_{q\mid x}(1-\frac1q)$
(произведение по простым делителям), то
\[
\frac{\varphi(L)}{\varphi(n)}=t\prod_{q\mid L,\ q\nmid n}\Bigl(1-\frac1q\Bigr)\ \ge\ t\prod_{q\mid t}\Bigl(1-\frac1q\Bigr)=\varphi(t)\ge1
\]
(простые $q\mid L$, не делящие $n$, делят $t$). Равенство $\varphi(L)=\varphi(n)$ требует $\varphi(t)=1$, т.\,е.\ $t\in\{1,2\}$,
а при $t=2$ ещё и $2\nmid n$ (иначе отношение равно $2$). Итак, $L=n$, или $n$ нечётно и $L=2n$; поэтому
$\xi\in\mu_L\subset\mu_n$ или $\mu_{2n}$. Обратно, $\mu_n\subset\Q(\zeta_n)$, а при нечётном $n$ и $\mu_{2n}\subset\Q(\zeta_n)$
по (a).

(c) Если $\Q(\zeta_m)\subset\Q(\zeta_n)$, то $\zeta_m\in\mu_n$ или $\mu_{2n}$ (при нечётном $n$) по (b). Обратно, при $m\mid n$
имеем $\zeta_m=\zeta_n^{n/m}$, а при нечётном $n$ и $m\mid2n$~--- $\zeta_m\in\Q(\zeta_{2n})=\Q(\zeta_n)$.
\end{proof}

\begin{proposition}[5.11c: композит и пересечение]
Пусть $d=\gcd(m,n)$ и $L=\operatorname{lcm}(m,n)$. Тогда
\[
\Q(\zeta_m)\,\Q(\zeta_n)=\Q(\zeta_L),\qquad \Q(\zeta_m)\cap\Q(\zeta_n)=\Q(\zeta_d).
\]
В частности, при $\gcd(m,n)=1$ имеем $\Q(\zeta_m)\cap\Q(\zeta_n)=\Q$ и $\Gal(\Q(\zeta_{mn})/\Q)\cong
\Gal(\Q(\zeta_m)/\Q)\times\Gal(\Q(\zeta_n)/\Q)$, что отвечает изоморфизму $(\Z/mn\Z)^\times\cong(\Z/m\Z)^\times\times(\Z/n\Z)^\times$.
\end{proposition}

\begin{proof}
Подгруппа $\langle\zeta_m,\zeta_n\rangle\subset\CC^\times$ конечна и потому циклическая (\rf{5.7c}); её порядок делится на $m$ и
$n$, т.\,е.\ на $L$, а сама она содержится в $\mu_L$. Значит, она равна $\mu_L$, и $\Q(\zeta_m,\zeta_n)$ содержит $\zeta_L$; обратно,
$\zeta_m,\zeta_n\in\Q(\zeta_L)$. Это доказывает первое равенство.

Пусть $K=\Q(\zeta_m)\cap\Q(\zeta_n)$. Расширение $\Q(\zeta_m)/\Q$ Галуа, поэтому по \rf{3.19}
$[\Q(\zeta_m)\Q(\zeta_n):\Q(\zeta_n)]=[\Q(\zeta_m):K]$, т.\,е.\ $\varphi(L)/\varphi(n)=\varphi(m)/[K:\Q]$ и
\[
[K:\Q]=\frac{\varphi(m)\varphi(n)}{\varphi(L)}=\varphi(d).
\]
Последнее равенство $\varphi(m)\varphi(n)=\varphi(L)\varphi(d)$ следует из формулы $\varphi(x)=x\prod_{q\mid x}(1-\frac1q)$,
равенства $mn=Ld$ и того, что простое $q$ делит $L$, если оно делит $m$ или $n$, и делит $d$, если делит оба. С другой
стороны, $\zeta_d=\zeta_m^{m/d}=\zeta_n^{n/d}\in K$ и $[\Q(\zeta_d):\Q]=\varphi(d)$. Значит, $K=\Q(\zeta_d)$. При $\gcd(m,n)=1$
отображение ограничения $\Gal(\Q(\zeta_{mn})/\Q)\to\Gal(\Q(\zeta_m)/\Q)\times\Gal(\Q(\zeta_n)/\Q)$ инъективно (автоморфизм,
тождественный на $\zeta_m$ и $\zeta_n$, тождествен на $\Q(\zeta_m,\zeta_n)$), а порядки групп равны $\varphi(mn)=\varphi(m)\varphi(n)$.
\end{proof}

Формулы верны при всех $m,n$, но $m$ и $n$ не восстанавливаются по полям (\rf{5.11b}(c)). Например,
$\Q(\zeta_6)\cap\Q(\zeta_{10})=\Q(\zeta_2)=\Q$ (т.\,е.\ $\Q(\sqrt{-3})\cap\Q(\zeta_5)=\Q$), а
$\Q(\zeta_6)\Q(\zeta_{10})=\Q(\zeta_{30})=\Q(\zeta_{15})$; и $\Q(\zeta_3)\cap\Q(\zeta_6)=\Q(\zeta_3)=\Q(\zeta_6)$.

\begin{proposition}[5.11d: максимальное вещественное подполе]
Пусть $n\ge3$ и $\zeta=\zeta_n$. Тогда
\[
\Q(\zeta)\cap\R=\Q(\zeta+\zeta^{-1})=\Q\bigl(\cos\tfrac{2\pi}n\bigr),
\]
это неподвижное поле комплексного сопряжения $\sigma_{-1}$, и $[\Q(\zeta+\zeta^{-1}):\Q]=\varphi(n)/2$. Подполе
$\Q(\zeta)^H$ вещественно тогда и только тогда, когда $-1\in H$.
\end{proposition}

\begin{proof}
Комплексное сопряжение переводит $\zeta$ в $\bar\zeta=\zeta^{-1}\in\Q(\zeta)$ и потому ограничивается до автоморфизма
$\sigma_{-1}$ поля $\Q(\zeta)$; при $n\ge3$ он нетривиален ($\zeta\notin\R$). Его неподвижное поле~--- $\Q(\zeta)\cap\R$,
степени $\varphi(n)/2$ над $\Q$. Число $\zeta+\zeta^{-1}=2\cos\frac{2\pi}n$ вещественно, так что
$\Q(\zeta+\zeta^{-1})\subset\Q(\zeta)\cap\R$. С другой стороны, $\zeta$~--- корень многочлена
\[
X^2-(\zeta+\zeta^{-1})X+1
\]
над $\Q(\zeta+\zeta^{-1})$, так что $[\Q(\zeta):\Q(\zeta+\zeta^{-1})]\le2=[\Q(\zeta):\Q(\zeta)\cap\R]$; отсюда равенство полей.
Наконец, $\Q(\zeta)^H\subset\R$ $\iff$ $\Q(\zeta)^H\subset\Q(\zeta)^{\{\pm1\}}$ $\iff$ $\{\pm1\}\subset H$.
\end{proof}

Например, минимальные многочлены $2\cos\frac{2\pi}n$: $X^2+X-1$ ($n=5$), $X^2-2$ ($n=8$), $X^2-3$ ($n=12$),
$X^3+X^2-2X-1$ ($n=7$, ср.\ \rf{3.22}), $X^3-3X+1$ ($n=9$), $X^5+X^4-4X^3-3X^2+3X+1$ ($n=11$),
$X^6+X^5-5X^4-4X^3+6X^2+3X-1$ ($n=13$); их степени равны $\varphi(n)/2$.

\begin{proposition}[5.11e: гауссовы периоды]
Пусть $p$~--- простое число, $\zeta=\zeta_p$ и $G=(\Z/p\Z)^\times$.
\begin{enumerate}[label=\textup{(\alph*)}]
\item Числа $\zeta,\zeta^2,\dots,\zeta^{p-1}$ образуют базис $\Q(\zeta)$ над $\Q$, и $G$ их переставляет:
$\sigma_b(\zeta^a)=\zeta^{ab}$.
\item Пусть $H\subset G$~--- подгруппа. \emph{Гауссов период} $\eta_H=\sum_{a\in H}\zeta^a$ порождает поле $\Q(\zeta)^H$; его
сопряжённые~--- это $[G:H]$ различных чисел $\eta_{bH}=\sum_{a\in bH}\zeta^a$ (по смежным классам $bH$), и минимальный
многочлен $\eta_H$ над $\Q$ равен $\prod_{bH\in G/H}(X-\eta_{bH})$.
\end{enumerate}
\end{proposition}

\begin{proof}
(a) Так как $\Phi_p$ неприводим, $1,\zeta,\dots,\zeta^{p-2}$~--- базис; из $1=-(\zeta+\dots+\zeta^{p-1})$ видно, что
$p-1$ чисел $\zeta,\dots,\zeta^{p-1}$ порождают $\Q(\zeta)$, а значит, тоже образуют базис.
(b) $\sigma_b(\eta_H)=\sum_{a\in H}\zeta^{ba}=\eta_{bH}$. Если $bH\ne H$, то $\eta_{bH}$ и $\eta_H$~--- суммы по
\emph{непересекающимся} наборам базисных векторов, и $\eta_{bH}\ne\eta_H$. Значит, стабилизатор $\eta_H$ в $G$ равен в
точности $H$, т.\,е.\ $\Gal(\Q(\zeta)/\Q(\eta_H))=H$ и $\Q(\eta_H)=\Q(\zeta)^H$. Сопряжённые $\eta_H$~--- это его образы
$\eta_{bH}$ под действием $G$, их $[G:H]$, откуда минимальный многочлен.
\end{proof}

Для составного $n$ это может нарушаться: для подгруппы $H=\{1,4,7\}\subset(\Z/9\Z)^\times$ (ядра редукции
$(\Z/9\Z)^\times\to(\Z/3\Z)^\times$) имеем $\Q(\zeta_9)^H=\Q(\zeta_3)$, но $\zeta_9+\zeta_9^4+\zeta_9^7=\zeta_9(1+\zeta_3+\zeta_3^2)=0$.

\begin{theorem}[5.11f: квадратичное подполе $\Q(\zeta_p)$ и сумма Гаусса]
Пусть $p$~--- нечётное простое число, $\zeta=\zeta_p$, $p^*=(-1)^{(p-1)/2}p$, и $\bigl(\frac ap\bigr)$~--- символ Лежандра
($=1$, если $a$~--- квадрат по модулю $p$, и $-1$ иначе; $p\nmid a$). Тогда $\Q(\zeta)$ содержит ровно одно квадратичное
подполе, и это $\Q(\sqrt{p^*})$. Точнее, \emph{сумма Гаусса}
\[
g=\sum_{a=1}^{p-1}\Bigl(\frac ap\Bigr)\zeta^a
\]
удовлетворяет $g^2=p^*$, $\sigma_b(g)=\bigl(\frac bp\bigr)g$, и квадратичное подполе равно $\Q(g)$.
\end{theorem}

\begin{proof}
Группа $G=(\Z/p\Z)^\times$ циклическая (\rf{5.7c}) чётного порядка $p-1$, поэтому в ней ровно одна подгруппа индекса $2$
(\rf{5.7b}(b))~--- подгруппа $R$ квадратов (отображение $x\mapsto x^2$ имеет ядро $\{\pm1\}$, так что квадратов $(p-1)/2$). По
\rf{5.11a} квадратичное подполе ровно одно~--- $\Q(\zeta)^R$. Символ Лежандра~--- гомоморфизм $G\to\{\pm1\}$ с ядром $R$
(изоморфизм $G/R\cong\{\pm1\}$).

\emph{Действие $G$.} $\sigma_b(g)=\sum_a\bigl(\frac ap\bigr)\zeta^{ab}=\sum_c\bigl(\frac{b^{-1}c}p\bigr)\zeta^c=\bigl(\frac bp\bigr)g$
(подстановка $c=ab$; $(\frac{b^{-1}}p)=(\frac bp)$).

\emph{Вычисление $g^2$.} Подставим в $g^2=\sum_{a,b}\bigl(\frac{ab}p\bigr)\zeta^{a+b}$ вместо $b$ выражение $at$: при
фиксированном $a$ элемент $t$ пробегает $G$, и $\bigl(\frac{a\cdot at}p\bigr)=\bigl(\frac tp\bigr)$. Получаем
\[
g^2=\sum_{t\in G}\Bigl(\frac tp\Bigr)\sum_{a\in G}\zeta^{a(1+t)} .
\]
Внутренняя сумма равна $p-1$ при $t=-1$ и $-1$ при $t\ne-1$ (тогда $a(1+t)$ пробегает все ненулевые классы, а
$\zeta+\dots+\zeta^{p-1}=-1$). Так как $\sum_{t\in G}\bigl(\frac tp\bigr)=0$ (квадратов и неквадратов поровну),
\[
g^2=\Bigl(\frac{-1}p\Bigr)(p-1)-\sum_{t\ne-1}\Bigl(\frac tp\Bigr)=\Bigl(\frac{-1}p\Bigr)(p-1)+\Bigl(\frac{-1}p\Bigr)=\Bigl(\frac{-1}p\Bigr)p .
\]
Наконец, $-1$~--- квадрат по модулю $p$ $\iff$ в циклической группе порядка $p-1$ есть элемент порядка $4$ (его квадрат~---
единственный элемент порядка $2$, т.\,е.\ $-1$) $\iff$ $4\mid p-1$; поэтому $\bigl(\frac{-1}p\bigr)=(-1)^{(p-1)/2}$ и $g^2=p^*$.

Итак, $g\ne0$, и $g$ неподвижен ровно относительно $R$; значит, $\Q(g)=\Q(\zeta)^R$ (основная теорема), а $g=\pm\sqrt{p^*}$.
\end{proof}

Если $\eta_0=\sum_{a\in R}\zeta^a$ и $\eta_1=\sum_{a\notin R}\zeta^a$~--- два гауссовых периода подгруппы $R$, то
$\eta_0+\eta_1=-1$ и $\eta_0-\eta_1=g$, так что $\eta_0=\frac{-1+g}2$. Алгебраически знак в равенстве
$g=\pm\sqrt{p^*}$ не определяется: $\sigma_b$ при неквадратном $b$ меняет $g$ на $-g$. При $\zeta=e^{2\pi i/p}$ он таков:
$g=\sqrt p$ при $p\equiv1\pmod4$ и $g=i\sqrt p$ при $p\equiv3\pmod4$ (теорема Гаусса; приводим её без доказательства~---
для конкретного $p$ знак проверяется численно). Примеры при $\zeta=e^{2\pi i/p}$: $p=3$: $\Q(\zeta_3)=\Q(\sqrt{-3})$; $p=5$:
$\eta_0=\zeta+\zeta^4=2\cos\frac{2\pi}5=\frac{-1+\sqrt5}2$, корень $X^2+X-1$; $p=7$: $\eta_0=\zeta+\zeta^2+\zeta^4=\frac{-1+\sqrt{-7}}2$
($\sqrt{-7}=i\sqrt7$), корень $X^2+X+2$; $p=13$: $\eta_0=\frac{-1+\sqrt{13}}2$, корень $X^2+X-3$.

Отсюда же следует, что \emph{каждое квадратичное поле $\Q(\sqrt d)$ ($d\ne0,1$ свободно от квадратов) содержится в
круговом поле}, например в $\Q(\zeta_{4|d|})$: число $\sqrt d$ с точностью до знака есть произведение нескольких из чисел
$\sqrt{-1}=\zeta_4$, $\sqrt2=\zeta_8+\zeta_8^{-1}$ и $\sqrt{p^*}\in\Q(\zeta_p)$ ($p\mid d$ нечётно), а все они лежат в
$\Q(\zeta_{4|d|})$ по \rf{5.11b}(c).

\subsection*{Примеры: решётки подполей}

\begin{example}[5.11g: $\Q(\zeta_8)$]
Здесь $\zeta=\zeta_8=\frac{1+i}{\sqrt2}$, $\Phi_8=X^4+1$, и $G=(\Z/8\Z)^\times=\{1,3,5,7\}$~--- четверная группа (квадрат каждого
элемента равен $1$); в частности, она не циклическая. Её подгруппы порядка $2$: $\{1,7\}$, $\{1,5\}$, $\{1,3\}$. Неподвижные
поля:
\begin{itemize}[leftmargin=*]
\item $\{1,7\}$ (комплексное сопряжение): $\zeta+\zeta^7=2\cos\frac\pi4=\sqrt2$, поле $\Q(\sqrt2)$ (\rf{5.11d});
\item $\{1,5\}$: $\sigma_5(\zeta)=\zeta^5=-\zeta$, а $\zeta^2=i$ неподвижен; поле $\Q(i)=\Q(\zeta_4)$;
\item $\{1,3\}$: $\sigma_3$ переставляет $\zeta$ и $\zeta^3$ ($\zeta^9=\zeta$), так что $\zeta+\zeta^3=i\sqrt2$ неподвижен; поле
$\Q(\sqrt{-2})$.
\end{itemize}
Значит, $\Q(\zeta_8)=\Q(i,\sqrt2)$; решётка подполей~--- на рисунке ниже слева.
\end{example}

\begin{example}[5.11h: $\Q(\zeta_7)$]
Пусть $\zeta=\zeta_7$. Группа $G=(\Z/7\Z)^\times$ циклическая порядка $6$ с образующей $3$ (степени $3$ по модулю $7$: $3,2,6,4,5,1$). Её подгруппы:
$\{1\}$, $\{1,6\}=\{\pm1\}$, $\{1,2,4\}$ (квадраты) и $G$. Неподвижные поля: $\Q(\zeta_7)^{\{\pm1\}}=\Q(\zeta+\zeta^{-1})$
степени $3$, порождённое $2\cos\frac{2\pi}7$ (корень $X^3+X^2-2X-1$), и $\Q(\zeta_7)^{\{1,2,4\}}=\Q(\sqrt{-7})$, порождённое
периодом $\zeta+\zeta^2+\zeta^4$ (корень $X^2+X+2$). Ср.\ \rf{3.22}.
\end{example}

\begin{keepfig}
\begin{center}
\begin{tikzpicture}[yscale=1.25,every node/.style={font=\small}]
\begin{scope}
\node (E) at (0,2.8) {$\Q(\zeta_8)=\Q(i,\sqrt2)$};
\node[cnew,align=center] (A) at (-2.1,1.4) {$\Q(\sqrt2)$\\[-1pt]{\scriptsize\color{cold}$\{1,7\}$}};
\node[cnew,align=center] (B) at (0,1.4) {$\Q(i)$\\[-1pt]{\scriptsize\color{cold}$\{1,5\}$}};
\node[cnew,align=center] (C) at (2.1,1.4) {$\Q(\sqrt{-2})$\\[-1pt]{\scriptsize\color{cold}$\{1,3\}$}};
\node (Q) at (0,0) {$\Q$};
\begin{scope}[every node/.style={font=\scriptsize}]
\draw (E) -- node[above left]{$2$} (A); \draw (E) -- node[right]{$2$} (B); \draw (E) -- node[above right]{$2$} (C);
\draw (A) -- node[below left]{$2$} (Q); \draw (B) -- node[right]{$2$} (Q); \draw (C) -- node[below right]{$2$} (Q);
\end{scope}
\end{scope}
\begin{scope}[xshift=7.8cm]
\node (E) at (0,2.8) {$\Q(\zeta_7)$};
\node[cnew,align=center] (A) at (-1.8,1.65) {$\Q(\zeta_7+\zeta_7^{-1})$\\[-1pt]{\scriptsize\color{cold}$\{\pm1\}$}};
\node[cnew,align=center] (B) at (1.8,1.05) {$\Q(\sqrt{-7})$\\[-1pt]{\scriptsize\color{cold}$\{1,2,4\}$}};
\node (Q) at (0,0) {$\Q$};
\begin{scope}[every node/.style={font=\scriptsize}]
\draw (E) -- node[above left]{$2$} (A); \draw (A) -- node[below left]{$3$} (Q);
\draw (E) -- node[above right]{$3$} (B); \draw (B) -- node[below right]{$2$} (Q);
\end{scope}
\end{scope}
\end{tikzpicture}
\end{center}
\figcap{Решётки подполей $\Q(\zeta_8)$ и $\Q(\zeta_7)$; на рёбрах~--- степени, под полями (синим)~--- соответствующие
подгруппы $H\subset(\Z/n\Z)^\times$, неподвижные поля которых изображены.}
\end{keepfig}

\begin{example}[5.11i: $\Q(\zeta_{13})$]
Пусть $\zeta=\zeta_{13}$. Группа $G=(\Z/13\Z)^\times$ циклическая порядка $12$ с образующей $2$ (степени $2$ по модулю $13$:
$2,4,8,3,6,12,11,9,5,10,7,1$); положим $\sigma=\sigma_2$. Подгруппы~--- $\langle\sigma^d\rangle$ для $d\mid12$ порядка $12/d$, а
их неподвижные поля $K_d=\Q(\zeta)^{\langle\sigma^d\rangle}$ имеют степень $d$, причём $K_d\subset K_{d'}\iff d\mid d'$. По
5.11e каждое $K_d$ порождается гауссовым периодом:
\begin{center}\small
\renewcommand{\arraystretch}{1.2}
\begin{tabular}{@{}clll@{}}
\toprule
$d$ & подгруппа $\langle 2^d\rangle\subset(\Z/13\Z)^\times$ & образующая $K_d$ & минимальный многочлен\\
\midrule
$2$ & $\{1,3,4,9,10,12\}$ (квадраты) & $\eta=\frac{-1+\sqrt{13}}2$ & $X^2+X-3$\\
$3$ & $\{1,5,8,12\}$ & $\alpha=\zeta+\zeta^5+\zeta^8+\zeta^{12}$ & $X^3+X^2-4X+1$\\
$4$ & $\{1,3,9\}$ & $\theta=\zeta+\zeta^3+\zeta^9$ & $X^4+X^3+2X^2-4X+3$\\
$6$ & $\{1,12\}=\{\pm1\}$ & $\zeta+\zeta^{-1}=2\cos\frac{2\pi}{13}$ & $X^6+X^5-5X^4-4X^3+6X^2+3X-1$\\
$12$ & $\{1\}$ & $\zeta$ & $\Phi_{13}$\\
\bottomrule
\end{tabular}
\end{center}
Поле $K_d$ вещественно $\iff$ $-1=2^6\in\langle2^d\rangle$ $\iff$ $d\mid6$ (5.11d): поля $K_2,K_3,K_6$ вещественны, а $K_4$ и
$K_{12}$~--- нет; например,
\[
\theta=\zeta+\zeta^3+\zeta^9=\frac{-1+\sqrt{13}+i\sqrt{26-6\sqrt{13}}}4 .
\]
Кроме того, $K_6=K_2K_3$ и $K_2=K_4\cap K_6$.
\end{example}

\begin{keepfig}
\begin{center}
\begin{tikzpicture}[yscale=1.35,every node/.style={font=\small}]
\begin{scope}
\node (e)  at (0,3)   {$\{1\}$};
\node (s4) at (-1.6,2) {$\langle\sigma^4\rangle$};
\node (s6) at (1.6,2)  {$\langle\sigma^6\rangle=\{\pm1\}$};
\node (s2) at (-1.6,1) {$\langle\sigma^2\rangle$};
\node (s3) at (1.6,1)  {$\langle\sigma^3\rangle$};
\node (G)  at (0,0)    {$G$};
\draw (e)--(s4) (e)--(s6) (s4)--(s2) (s6)--(s2) (s6)--(s3) (s2)--(G) (s3)--(G);
\node[cold,font=\scriptsize] at (0,-0.55) {подгруппы (вверху~--- меньшие)};
\end{scope}
\node at (3.9,1.5) {$\longleftrightarrow$};
\begin{scope}[xshift=7.6cm]
\node (K12) at (0,3)   {$\Q(\zeta_{13})$};
\node[cnew] (K4)  at (-1.7,2) {$K_4=\Q(\theta)$};
\node[cnew] (K6)  at (1.7,2)  {$K_6=\Q(\cos\frac{2\pi}{13})$};
\node[cnew] (K2)  at (-1.7,1) {$K_2=\Q(\sqrt{13})$};
\node[cnew] (K3)  at (1.7,1)  {$K_3=\Q(\alpha)$};
\node (K1)  at (0,0)    {$\Q$};
\begin{scope}[every node/.style={font=\scriptsize}]
\draw (K12)--node[above left]{3}(K4) (K12)--node[above right]{2}(K6)
      (K4)--node[left]{2}(K2) (K6)--node[above left]{3}(K2)
      (K6)--node[right]{2}(K3) (K2)--node[below left]{2}(K1) (K3)--node[below right]{3}(K1);
\end{scope}
\node[cold,font=\scriptsize] at (0,-0.55) {подполя (вверху~--- бóльшие)};
\end{scope}
\end{tikzpicture}
\end{center}
\figcap{Решётка подгрупп циклической группы $G=\langle\sigma\rangle\cong\Z/12$ и соответствующая ей (по основной теореме)
решётка подполей $\Q(\zeta_{13})$; на рёбрах справа~--- степени.}
\end{keepfig}

\subsection*{Элементы Фробениуса и их применения}

Пусть $p\nmid n$~--- простое число. Положим $\sigma_p=\sigma_{[p]}\in\Gal(\Q(\zeta_n)/\Q)$, т.\,е.\ $\sigma_p(\zeta)=\zeta^p$ для
$\zeta\in\mu_n$. Второе доказательство теоремы \rf{5.10} показывает: \emph{для любого максимального идеала $P$ кольца
$\Z[\zeta_n]$, содержащего $p$,}
\[
\sigma_p(a)\equiv a^p\pmod P\qquad\text{для всех } a\in\Z[\zeta_n],
\]
т.\,е.\ $\sigma_p$~--- элемент Фробениуса в каждом $P$ над $p$ (там мы показали, что $\sigma_P(\zeta)=\zeta^p$, т.\,е.\
$\sigma_P=\sigma_p$). Поскольку группа Галуа абелева, <<класс Фробениуса>> в $p$ (\rf{4.29}) состоит из одного элемента $\sigma_p$.
Кроме доказательства \rf{5.10}, элементы Фробениуса дают следующее.

\begin{proposition}[5.11j: разложение $\Phi_n$ по модулю $p$]
Пусть $p\nmid n$ и $f$~--- порядок $[p]$ в $(\Z/n\Z)^\times$. Тогда $\Phi_n$ по модулю $p$ раскладывается в произведение
$\varphi(n)/f$ различных унитарных неприводимых многочленов степени $f$ из $\F_p[X]$. Если $\xi$~--- первообразный корень
$n$-й степени из $1$ в $\Fpb$, то корни неприводимого множителя, делящегося на $X-\xi^a$,~--- это $\xi^{b}$, $b\in a\langle p\rangle$:
множители отвечают смежным классам $(\Z/n\Z)^\times/\langle[p]\rangle$.
\end{proposition}

\begin{proof}
По \rf{5.8d}(b) корни $\bar\Phi_n$ в $\Fpb$~--- это $\varphi(n)$ первообразных корней $n$-й степени из $1$, все простые. Пусть
$\xi$~--- один из них. Группа $\Gal(\F_p(\xi)/\F_p)$ порождается автоморфизмом $x\mapsto x^p$ (\rf{4.20}), поэтому сопряжённые с $\xi$
над $\F_p$~--- это $\xi,\xi^p,\xi^{p^2},\dots$, а $\xi^{p^k}=\xi\iff n\mid p^k-1\iff f\mid k$. Значит, минимальный многочлен $\xi$
имеет ровно $f$ корней $\xi,\xi^p,\dots,\xi^{p^{f-1}}$, т.\,е.\ степень $f$. Каждый неприводимый множитель $\bar\Phi_n$~---
минимальный многочлен любого своего корня; все они имеют степень $f$ и попарно различны (корни простые), и их
$\varphi(n)/f$. Корни множителя, содержащего $\xi^a$,~--- это $\xi^{ap^k}$: показатели умножаются на степени $p$.
\end{proof}

На языке группы Галуа: $\sigma_p$ имеет в $\Gal(\Q(\zeta_n)/\Q)\cong(\Z/n\Z)^\times$ порядок $f$ и переставляет $\varphi(n)$ корней
$\Phi_n$ циклами длины $f$~--- в точности как предсказывает теорема Дедекинда \rf{4.28} по разложению $\Phi_n$ по модулю $p$.
Примеры:
\begin{itemize}[leftmargin=*]
\item $n=7$, $p=2$: $\langle2\rangle=\{1,2,4\}$, смежные классы $\{1,2,4\}$ и $\{3,5,6\}$, и $\Phi_7\equiv(X^3+X+1)(X^3+X^2+1)\pmod2$.
\item $n=5$: при $p\equiv1\pmod5$ четыре линейных множителя (при $p=11$ корни $3,4,5,9$); при $p\equiv4\pmod5$ два
квадратичных (при $p=19$: $(X^2+5X+1)(X^2-4X+1)$); при $p\equiv2,3\pmod5$ многочлен $\Phi_5$ неприводим (например, $p=2,3$).
\item $n=8$: в $(\Z/8\Z)^\times$ все элементы имеют порядок $\le2$, поэтому $\Phi_8=X^4+1$, неприводимый над $\Q$, приводим по
модулю \emph{каждого} простого: по модулю $2$ он равен $(X+1)^4$, а при нечётном $p$ раскладывается на множители степени
$\le2$; например, $X^4+1\equiv(X^2+X+2)(X^2+2X+2)\pmod3$.
\end{itemize}
В частности, при $p\nmid n$ многочлен $\Phi_n$ неприводим по модулю $p$ тогда и только тогда, когда $[p]$ порождает
$(\Z/n\Z)^\times$ (так что при нецикличной $(\Z/n\Z)^\times$ таких $p$ нет), и раскладывается на линейные множители тогда и
только тогда, когда $p\equiv1\pmod n$. Условие $p\nmid n$ здесь существенно: например, $\Phi_6=X^2-X+1\equiv X^2+X+1\pmod2$
неприводим над $\F_2$, хотя $[2]\notin(\Z/6\Z)^\times$.

\begin{corollary}[5.11k]
Для каждого $n\ge1$ существует бесконечно много простых чисел $p\equiv1\pmod n$.
\end{corollary}

\begin{proof}
\emph{Если $a\in\Z$, $q$~--- простое, $q\mid\Phi_n(a)$ и $q\nmid n$, то $q\equiv1\pmod n$.} Действительно, $\bar a\in\F_q$~---
корень $\bar\Phi_n$, т.\,е.\ по \rf{5.8d}(b) первообразный корень $n$-й степени из $1$ в $\F_q$; его порядок $n$ делит
$|\F_q^\times|=q-1$.

Подставив $X=0$ в \eqref{eq:W1-Xn-1}, получаем $-1=\prod_{d\mid n}\Phi_d(0)$, так что все $\Phi_d(0)=\pm1$. Пусть
$p_1,\dots,p_r$~--- уже найденные простые $\equiv1\pmod n$, и $a=Nnp_1\cdots p_r$, где $N$ столь велико, что $|\Phi_n(a)|>1$
(многочлен $\Phi_n$ непостоянен). Так как $\Phi_n(a)\equiv\Phi_n(0)=\pm1\pmod a$, любой простой делитель $q$ числа
$\Phi_n(a)$ не делит $a$; значит, $q\nmid n$, $q\ne p_i$, и по доказанному $q\equiv1\pmod n$~--- новое простое.
\end{proof}

Следующее применение элемента Фробениуса~--- закон взаимности для символа Лежандра. Нам понадобится \emph{критерий
Эйлера}: для нечётного простого $q$ и $q\nmid a$
\[
a^{(q-1)/2}\equiv\Bigl(\frac aq\Bigr)\pmod q .
\]
(Отображение $x\mapsto x^{(q-1)/2}$~--- гомоморфизм циклической группы $\F_q^\times$ в $\{\pm1\}$; его ядро~--- корни
$X^{(q-1)/2}-1$, их не больше $(q-1)/2$, и все $(q-1)/2$ квадратов в нём лежат; значит, ядро~--- в точности квадраты.)

\begin{theorem}[5.11l: квадратичный закон взаимности]
Для различных нечётных простых $p$ и $q$
\[
\Bigl(\frac qp\Bigr)\Bigl(\frac pq\Bigr)=(-1)^{\frac{p-1}2\cdot\frac{q-1}2}.
\]
\end{theorem}

\begin{proof}
Пусть $\zeta=\zeta_p$, $A=\Z[\zeta]$, $g$~--- сумма Гаусса (\rf{5.11f}) и $P$~--- максимальный идеал $A$, содержащий $q$ (он
существует по той же причине, что и во втором доказательстве \rf{5.10}); тогда $P\cap\Z=q\Z$. Так как $g\in A$, то по свойству элемента Фробениуса $\sigma_q(g)\equiv g^q\pmod P$. Левая часть равна
$\bigl(\frac qp\bigr)g$ (\rf{5.11f}), а правая~--- $g\,(g^2)^{(q-1)/2}=g\,(p^*)^{(q-1)/2}$. Умножив сравнение
$\bigl(\frac qp\bigr)g\equiv(p^*)^{(q-1)/2}g$ на $g$ и использовав $g^2=p^*$, получаем сравнение целых чисел
\[
\Bigl(\frac qp\Bigr)p^*\equiv(p^*)^{(q+1)/2}\pmod P,\quad\text{т.\,е.}\quad \pmod{P\cap\Z=q\Z}.
\]
Сократив на $p^*$ (оно взаимно просто с $q$) и применив критерий Эйлера, получаем
$\bigl(\frac qp\bigr)\equiv(p^*)^{(q-1)/2}\equiv\bigl(\frac{p^*}q\bigr)\pmod q$; обе части равны $\pm1$ и $q>2$, поэтому
$\bigl(\frac qp\bigr)=\bigl(\frac{p^*}q\bigr)$. Наконец, символ Лежандра мультипликативен, и
$\bigl(\frac{-1}q\bigr)=(-1)^{(q-1)/2}$ (критерий Эйлера), так что
$\bigl(\frac{p^*}q\bigr)=(-1)^{\frac{p-1}2\cdot\frac{q-1}2}\bigl(\frac pq\bigr)$.
\end{proof}

Смысл: $\bigl(\frac{p^*}q\bigr)=1$ означает, что элемент Фробениуса $\sigma_q$ тривиален на $\Q(\sqrt{p^*})$, т.\,е.\ $[q]$
лежит в подгруппе квадратов $R$; а это и есть $\bigl(\frac qp\bigr)=1$.

\begin{aside}[5.11m: кольца целых и классы Фробениуса]
Пусть $E/\Q$~--- конечное расширение Галуа с группой $G$, $\mathcal O_E$~--- кольцо целых (элементы $E$, являющиеся корнями
унитарных многочленов из $\Z[X]$), и $P$~--- простой идеал $\mathcal O_E$ с $P\cap\Z=p\Z$. Если $p$ не разветвлено в $E$, то
существует единственный $\sigma_P\in G$ с $\sigma_PP=P$ и $\sigma_P(a)\equiv a^p\pmod P$ для всех $a\in\mathcal O_E$; при
замене $P$ на $\tau P$ он заменяется на $\tau\sigma_P\tau^{-1}$, и класс сопряжённости~--- \emph{класс Фробениуса} в $p$ (\rf{4.29};
в доказательстве \rf{4.28} вместо $\mathcal O_E$ использовалось кольцо, порождённое корнями многочлена). Это предмет
алгебраической теории чисел; приведём без доказательства, что даёт эта теория для $E=\Q(\zeta_n)$:
$\mathcal O_E=\Z[\zeta_n]$; простые $p\nmid n$ не разветвлены, и $p\mathcal O_E$~--- произведение $\varphi(n)/f$ различных
простых идеалов с полями вычетов $\F_{p^f}$, $f=\operatorname{ord}_n(p)$ (в согласии с \rf{5.11j}); класс Фробениуса равен
$\{\sigma_p\}$. По теореме Чеботарёва каждый класс сопряжённости $G$ является классом Фробениуса для бесконечного
множества простых; для $\Q(\zeta_n)$ это теорема Дирихле о простых в арифметических прогрессиях (\rf{5.11k}~--- её
элементарный частный случай). Наконец, по теореме Кронекера--Вебера каждое конечное абелево расширение $\Q$ содержится в
некотором круговом поле; для квадратичных полей мы проверили это выше.
\end{aside}

\subsection*{Построение правильных многоугольников}

\begin{theorem}[5.12: Гаусс]
Правильный $n$-угольник ($n\ge3$) можно построить циркулем и линейкой тогда и только тогда, когда
$n=2^kp_1\cdots p_s$, где $k\ge0$, а $p_i$~--- различные простые числа Ферма, т.\,е.\ простые вида $2^{2^r}+1$.
\end{theorem}

\begin{proof}
Правильный $n$-угольник построим тогда и только тогда, когда построимо число $\cos\frac{2\pi}n$: вершины многоугольника,
вписанного в единичную окружность,~--- точки $P_k=(\cos\frac{2\pi k}n,\sin\frac{2\pi k}n)$, $0\le k<n$; по точке $P_1$
строится её проекция $\cos\frac{2\pi}n$ на ось, а по $\cos\frac{2\pi}n$~--- перпендикуляр к оси, пересекающий окружность в
$P_1$, после чего сторона $P_0P_1$ откладывается по окружности.

Число $2\cos\frac{2\pi}n=\zeta_n+\zeta_n^{-1}$ порождает максимальное вещественное подполе $\Q(\zeta_n)\cap\R$ степени
$\varphi(n)/2$ (\rf{5.11d}); это расширение Галуа поля $\Q$ (\rf{5.11a}(b)). Если $\cos\frac{2\pi}n$ построимо, то
$[\Q(\cos\frac{2\pi}n):\Q]=\varphi(n)/2$~--- степень двойки (\rf{1.38}), и $\varphi(n)$~--- степень двойки. Обратно, если
$\varphi(n)=2^m$, то $\Q(\cos\frac{2\pi}n)\subset\R$~--- расширение Галуа поля $\Q$ степени $2^{m-1}$, и $\cos\frac{2\pi}n$
построимо по \rf{3.24}.

Осталось выяснить, когда $\varphi(n)$~--- степень двойки. Пусть $n=q_1^{e_1}\cdots q_t^{e_t}$~--- разложение на простые
множители ($q_i$ различны, $e_i\ge1$). Так как $\varphi$ мультипликативна (китайская теорема об остатках) и
$\varphi(q^e)=q^e-q^{e-1}=q^{e-1}(q-1)$,
\[
\varphi(n)=\prod_{i=1}^tq_i^{e_i-1}(q_i-1).
\]
Это степень двойки тогда и только тогда, когда для каждого $i$ степенями двойки являются и $q_i^{e_i-1}$, и $q_i-1$. Для
$q_i=2$ это выполнено при любом $e_i$ ($q_i-1=1$). Для нечётного $q_i$ нужно $e_i=1$ (иначе $q_i\mid\varphi(n)$) и
$q_i=2^w+1$. Наконец, если $2^w+1$ простое ($w\ge1$), то $w$~--- степень двойки: если $w=ab$ с нечётным $b>1$, то
$2^a+1$ делит $(2^a)^b+1=2^w+1$ (так как $x+1\mid x^b+1$ при нечётном $b$), причём $1<2^a+1<2^w+1$. Значит, каждый
нечётный $q_i$ входит в $n$ в первой степени и имеет вид $2^{2^r}+1$, т.\,е.\ является простым числом Ферма, и $n$ имеет
указанный вид. Обратно, если $n=2^kp_1\cdots p_s$ с различными простыми $p_i=2^{2^{r_i}}+1$, то число
$\varphi(n)=2^{k-1}\prod_i2^{2^{r_i}}$ (или $\prod_i2^{2^{r_i}}$ при $k=0$)~--- степень двойки.
\end{proof}

\begin{remark}[5.13]
Достаточность условия теоремы \rf{5.12} доказал Гаусс в <<Арифметических исследованиях>> (1801). Числа Ферма~--- это
$\mathrm F_r=2^{2^r}+1$. Числа $\mathrm F_0,\dots,\mathrm F_4=3,5,17,257,65537$ простые, а
$\mathrm F_5=4294967297=641\cdot6700417$ (Эйлер), и $\mathrm F_5,\dots,\mathrm F_{10}$ составные. Неизвестно, есть ли другие простые числа Ферма, поэтому список $n$, для которых правильный
$n$-угольник построим, до сих пор не известен полностью. Например, при $3\le n\le100$ это
\[
3,4,5,6,8,10,12,15,16,17,20,24,30,32,34,40,48,51,60,64,68,80,85,96 .
\]
Если $\mathrm F_0,\dots,\mathrm F_4$~--- все простые Ферма, то нечётных таких $n$ ровно $2^5-1=31$, наибольшее из них
$3\cdot5\cdot17\cdot257\cdot65537=2^{32}-1$. Самый знаменитый случай~--- $17$-угольник ($\varphi(17)=16$): цепочка подполей
$\Q\subset K_2\subset K_4\subset K_8\subset\Q(\zeta_{17})$ с шагами степени $2$ (подгруппы циклической группы порядка $16$)
даёт построение через последовательное решение квадратных уравнений.
\end{remark}

\begin{remark}[5.13a: семиугольник и одиннадцатиугольник]
Правильный $7$-угольник циркулем и линейкой не построим: $\varphi(7)=6$, и $2\cos\frac{2\pi}7$ имеет степень $3$ (минимальный
многочлен $X^3+X^2-2X-1$). Однако его можно построить, если разрешено ещё делить углы на три части (трисектор): корни
кубического многочлена с тремя вещественными корнями выражаются через косинус трети угла, и
\[
\cos\frac{2\pi}7=-\frac16+\frac{\sqrt7}3\cos\Bigl(\frac13\arccos\frac1{2\sqrt7}\Bigr)
\]
(обе части равны $0{,}6234898\ldots$; ср.\ упражнение 1-4). Правильный $11$-угольник нельзя построить и
с трисектором: $2\cos\frac{2\pi}{11}$ имеет степень $\varphi(11)/2=5$ (минимальный многочлен $X^5+X^4-4X^3-3X^2+3X+1$), а
построения циркулем, линейкой и трисектором дают лишь башни расширений степеней $2$ и $3$ (пересечения прямых и окружностей~---
уравнения степени $\le2$, деление угла $\theta$ на три~--- уравнение $4x^3-3x=\cos\theta$), так что степени построимых чисел
имеют вид $2^a3^b$.
\end{remark}
